\scanpage{48}\mspage{4}

\deleted{soit $(e_{\deleted{\unsure{a}}},p)$, avec}
\[
\deleted{p:e\longrightarrow a.}
\]

\deleted{Si $b$ est dans $X$, on sait \inserted{de $=$} que \inserted{dans} la
catégorie $X_{/b}$ il y a un plus petit
objet, et on voit \deleted{\unsure{tout de suite}} \inserted{facilement} qu'on
peut prendre sa source égale à $e$,
donc le prendre de la forme}
\[
\deleted{\deleted{\unsure{\longrightarrow b}}\qquad e\xrightarrow{p_b}b.}
\]

\deleted{Mais \deleted{\unsure{a pr}} il n'y a pas unicité \unsure{de|des}
$p_b$, et pour assurer une unicité, si
faire se peut, on voudrait mettre
en relation le choix de $p_b$, avec
celui (déjà fait) de $p_a=p$. Comme
$a$ est un plus grd objet de $X$,
il existe}
\[
\deleted{u:b\xrightarrow{\unsure{\circ}}a\,\unsure{x|r}},
\]
\deleted{et on voit facilement aussi que pour
$u$ fixé, on peut choisir $p_b$ de
la forme $p_{bu}$, avec}
\diagram{\begin{tikzcd}[ampersand replacement=\&,row sep=small,column sep=small,execute at end picture={\draw (current bounding box.south west)--(current bounding box.north east);}]
\& \deleted{b}\arrow[dr,"\deleted{u}"] \& {}\\
\deleted{e}\arrow[ur,"\deleted{p_u}"]\arrow[rr,"\deleted{p}"'] \& {} \& \deleted{a}
\end{tikzcd}}
\deleted{commutatif.}

% archive sheet number: 47
